\section{\textit{Support Vector Machine} (SVM)}
\label{sec:svm}

\textit{Support Vector Machine} (SVM) mencari \textit{hyperplane} pemisah dua kelas dengan margin terbesar, dan lewat \textit{kernel trick} dapat membentuk batas keputusan yang nonlinear tanpa membangun fitur nonlinear secara eksplisit \cite{ruskanda-svm-what,ruskanda-svm-linear,ruskanda-svm-nonlinear}. Model ini cocok untuk data pembatalan pesanan karena batas antara pesanan yang batal dan yang tidak batal tidak linear terhadap fitur (kernel Laplacian jauh mengungguli kernel linear, Subbab~\ref{subsec:svm-validasi}), jumlah barisnya (27.309) masih terjangkau oleh pelatihan berbasis kernel, dan model hanya ditentukan oleh \textit{support vector}, yaitu titik di sekitar dan di dalam margin. Konsekuensinya, SVM sensitif terhadap skala fitur dan tidak menerima nilai kosong, sehingga \textit{preprocessing} menentukan hasilnya (Subbab~\ref{subsec:svm-preprocessing}).

Implementasi \textit{from scratch} berada pada kelas \texttt{SVM} di \nolinkurl{src/algoritma/svm/svm.py} dan hanya memakai NumPy. Kelas ini mendukung kernel \textit{linear}, \textit{polynomial}, dan \textit{Gaussian} (RBF) yang diwajibkan, ditambah kernel Laplacian. Masalah dual diselesaikan dengan \textit{Sequential Minimal Optimization} (SMO) \cite{platt1998smo}, dengan pemilihan pasangan variabel berdasarkan informasi orde kedua \cite{fan2005wss} dan \textit{cache} baris kernel seperti pada LIBSVM \cite{chang2011libsvm}. Tiga kernel wajib divalidasi terhadap \texttt{sklearn.svm.SVC} \cite{pedregosa2011sklearn} dan hasilnya setara (Subbab~\ref{subsec:svm-hasil}).

Fokus bab ini adalah model terbaik, yaitu SVM dengan kernel Laplacian $K(a, b) = \exp(-\lVert W(a-b)\rVert_1 / 20)$, $C = 2$, dan 46 fitur yang distandardisasi lalu dikalikan bobot fitur $W$ yang dipelajari dari data latih. Fitur tersebut mencakup fitur ``saudara'', yaitu rata-rata label pesanan lain yang kemungkinan satu rombongan (Subbab~\ref{subsec:svm-preprocessing}). Dengan \textit{stratified} 5-\textit{fold cross validation} pada seluruh data latih bersih, model ini mencapai Macro F1 \textit{out-of-fold} 0,8703 (rata-rata 7 \textit{seed}, simpangan baku antar-\textit{seed} 0,0007) dan 0,8699 pada empat \textit{seed} konfirmasi yang tidak dipakai memilih konfigurasi. Skor publik Kaggle model yang dilatih pada seluruh data latih adalah 0,87464 (Bagian~\ref{sec:kaggle}). Sebagai pembanding, kernel wajib terbaik pada protokol notebook (hold-out 20\%, 6.000 baris latih) hanya mencapai Macro F1 0,8251, sedangkan kernel Laplacian dengan bobot fitur mencapai 0,8561 pada protokol yang sama. Keunggulan ini muncul dari tiga hal yang dibahas bertahap pada subbab berikut: jarak L1 pada kernel, bobot fitur, dan fitur saudara. Selisih SVM terhadap pohon keputusan C4.5 dengan 46 fitur yang sama (0,8685) kecil dan selang kepercayaannya masih mencakup nol (Subbab~\ref{subsec:svm-validasi}).

\subsection{Implementasi \textit{From Scratch}}
\label{subsec:svm-implementasi}

\subsubsection{Landasan Teori}

\paragraph{\textit{Hyperplane} dan margin.} Misalkan data latih $\{(x_i, y_i)\}_{i=1}^{n}$ dengan $x_i \in \mathbb{R}^d$ dan label $y_i \in \{-1, +1\}$. Pada implementasi, kelas batal diberi label $+1$ dan kelas tidak batal $-1$. SVM mencari \textit{hyperplane} $w \cdot x + b = 0$ yang memisahkan kedua kelas. Jarak titik $x$ ke \textit{hyperplane} adalah $|w \cdot x + b| / \lVert w\rVert$. Jika titik terdekat dari kedua kelas dinormalkan sehingga memenuhi $y_i(w \cdot x_i + b) = 1$, lebar margin antara kedua kelas adalah $2 / \lVert w\rVert$, sehingga memaksimumkan margin sama dengan meminimumkan $\lVert w\rVert$ \cite{ruskanda-svm-what,ruskanda-svm-linear}. Titik yang berada tepat di batas margin disebut \textit{support vector}.

\paragraph{\textit{Soft margin} dan parameter $C$.} Data nyata tidak dapat dipisahkan sempurna, sehingga setiap titik diberi variabel kendur (\textit{slack}) $\xi_i \ge 0$ yang mengukur seberapa jauh titik itu melanggar margin \cite{ruskanda-svm-nonlinear}. Masalah primalnya adalah
\begin{equation}
	\min_{w,\,b,\,\xi}\ \tfrac{1}{2}\lVert w\rVert^2 + \sum_{i=1}^{n} C_i\,\xi_i
	\quad \text{dengan} \quad y_i\big(w \cdot \varphi(x_i) + b\big) \ge 1 - \xi_i,\ \ \xi_i \ge 0,
	\label{eq:svm-primal}
\end{equation}
dengan $\varphi$ pemetaan fitur ke ruang berdimensi lebih tinggi (dipakai secara implisit lewat kernel). Pada implementasi $C_i = C$ untuk semua titik, kecuali jika \texttt{class\_weight} diberikan: $C_i = C \cdot c_{y_i}$ dengan $c_{y_i}$ bobot kelas titik tersebut (dipakai pada bonus \textit{cost-sensitive}, Subbab~\ref{subsubsec:svm-cost}). $C$ besar menghukum pelanggaran margin dengan keras sehingga margin menyempit dan model cenderung \textit{overfit}; $C$ kecil membiarkan lebih banyak pelanggaran sehingga margin lebar dan model cenderung \textit{underfit}. Pengaruh ini terlihat langsung pada hasil \textit{tuning} (Subbab~\ref{subsec:svm-validasi}).

\paragraph{Lagrangian dan bentuk dual.} Dengan pengali Lagrange $\alpha_i \ge 0$ untuk kendala margin dan $\mu_i \ge 0$ untuk $\xi_i \ge 0$,
\begin{equation}
	\mathcal{L}(w, b, \xi, \alpha, \mu) = \tfrac{1}{2}\lVert w\rVert^2 + \sum_i C_i\xi_i - \sum_i \alpha_i\big[y_i(w \cdot \varphi(x_i) + b) - 1 + \xi_i\big] - \sum_i \mu_i\xi_i .
	\label{eq:svm-lagrangian}
\end{equation}
Menurunkan $\mathcal{L}$ terhadap variabel primal dan menyamakannya dengan nol memberi
\begin{align}
	\frac{\partial \mathcal{L}}{\partial w} = 0 &\Rightarrow w = \sum_i \alpha_i y_i\,\varphi(x_i), \\
	\frac{\partial \mathcal{L}}{\partial b} = 0 &\Rightarrow \sum_i \alpha_i y_i = 0, \\
	\frac{\partial \mathcal{L}}{\partial \xi_i} = 0 &\Rightarrow \alpha_i = C_i - \mu_i .
\end{align}
Karena $\mu_i \ge 0$, diperoleh kotak batas $0 \le \alpha_i \le C_i$. Substitusi ke $\mathcal{L}$ menghilangkan $w$, $b$, dan $\xi$, sehingga tersisa masalah dual yang hanya bergantung pada $\alpha$ dan hasil kali dalam $K(x_i, x_j) = \varphi(x_i)\cdot\varphi(x_j)$:
\begin{equation}
	\max_{\alpha}\ \sum_{i} \alpha_i - \tfrac{1}{2}\sum_{i}\sum_{j} \alpha_i\alpha_j\,y_i y_j\,K(x_i, x_j)
	\quad \text{dengan} \quad 0 \le \alpha_i \le C_i,\ \ \sum_i \alpha_i y_i = 0 .
	\label{eq:svm-dual}
\end{equation}
Fungsi keputusan model hanya melibatkan titik dengan $\alpha_i > 0$ (\textit{support vector}):
\begin{equation}
	f(x) = \sum_{i} \alpha_i y_i\,K(x_i, x) + b, \qquad \hat{y} = \begin{cases} +1 & f(x) > t \\ -1 & \text{lainnya,} \end{cases}
	\label{eq:svm-decision}
\end{equation}
dengan $t$ ambang keputusan (\texttt{threshold} pada \texttt{predict}, bawaan 0). Nilai $f(x)$ juga disebut skor keputusan, dan $|f(x)|/\lVert w\rVert$ adalah jarak $x$ ke \textit{hyperplane} di ruang fitur kernel.

\paragraph{Kondisi KKT.} Solusi optimum memenuhi kondisi \textit{Karush-Kuhn-Tucker} berikut, yang menjadi dasar kriteria berhenti SMO:
\begin{equation}
	\alpha_i = 0 \Rightarrow y_i f(x_i) \ge 1, \qquad
	0 < \alpha_i < C_i \Rightarrow y_i f(x_i) = 1, \qquad
	\alpha_i = C_i \Rightarrow y_i f(x_i) \le 1 .
	\label{eq:svm-kkt}
\end{equation}
Titik dengan $\alpha_i = 0$ berada di luar margin dan tidak memengaruhi model. Titik dengan $0 < \alpha_i < C_i$ (\textit{free support vector}) tepat di batas margin. Titik dengan $\alpha_i = C_i$ berada di dalam margin atau salah sisi.

\paragraph{\textit{Kernel trick} dan empat kernel.} Persamaan \eqref{eq:svm-dual} dan \eqref{eq:svm-decision} hanya memakai $K(a, b) = \varphi(a)\cdot\varphi(b)$, sehingga $\varphi$ tidak pernah dihitung. Cukup mengganti $K$ untuk mendapatkan batas keputusan nonlinear \cite{ruskanda-svm-nonlinear}. Tabel~\ref{tab:svm-kernel} merangkum empat kernel yang diimplementasikan. Kernel \textit{sigmoid} pada materi kelas tidak diimplementasikan karena spesifikasi hanya mewajibkan tiga kernel pertama.

\begin{table}[H]
	\centering
	\caption{Kernel yang diimplementasikan pada \texttt{SVM}}
	\label{tab:svm-kernel}
	\small
	\begin{tabularx}{\textwidth}{l l Y}
		\toprule
		\textbf{\texttt{kernel}} & $K(a, b)$ & \textbf{Parameter dan catatan} \\
		\midrule
		\texttt{"linear"} & $a \cdot b$ & Tanpa parameter. Koefisien $w = \sum_i \alpha_i y_i x_i$ tersedia pada atribut \texttt{coef\_}. \\
		\texttt{"poly"} & $(\gamma\,a \cdot b + r)^{p}$ & $\gamma$ (\texttt{gamma}), $r$ (\texttt{coef0}), derajat $p$ (\texttt{degree}). \\
		\texttt{"rbf"} (\textit{Gaussian}) & $\exp\!\big(-\gamma\lVert a - b\rVert_2^2\big)$ & $\gamma$ (\texttt{gamma}). Jarak L2 kuadrat. \\
		\texttt{"laplacian"} & $\exp\!\big(-\gamma\lVert a - b\rVert_1\big)$ & $\gamma$ (\texttt{gamma}). Jarak L1 (jumlah selisih mutlak per fitur). \\
		\bottomrule
	\end{tabularx}
\end{table}

Nilai \texttt{gamma="scale"} (bawaan) berarti $\gamma = 1/(d \cdot \operatorname{Var}(X))$ dengan $\operatorname{Var}(X)$ variansi seluruh elemen matriks latih, dan \texttt{gamma="auto"} berarti $\gamma = 1/d$, sama seperti pada \textit{scikit-learn} \cite{pedregosa2011sklearn}. Pada data yang sudah distandardisasi, $\operatorname{Var}(X) \approx 1$ sehingga \texttt{"scale"} kira-kira $1/d$. Kernel \textit{Gaussian} dan Laplacian memenuhi $K(a, a) = 1$, sedangkan kernel \textit{linear} dan \textit{polynomial} tidak. Karena itu implementasi menghitung diagonal kernel secara terpisah (\texttt{\_gram\_diagonal} pada Tabel~\ref{tab:svm-fungsi}), tanpa membentuk seluruh matriks.

\paragraph{Kernel Laplacian dan alasan jarak L1.} Kernel Laplacian termasuk keluarga kernel $\exp(-\rho\sum_k |a_k - b_k|^{\beta})$ yang dipakai untuk klasifikasi histogram citra oleh \cite{chapelle1999histogram}, dengan $\beta = 1$. Perbedaannya dengan \textit{Gaussian} hanya pada jarak: Gaussian memakai jumlah kuadrat selisih, Laplacian memakai jumlah selisih mutlak. Pada data ini ada dua alasan jarak L1 lebih sesuai. Pertama, jarak L2 kuadrat memperbesar selisih ekstrem secara kuadratik. Setelah standardisasi, satu penanda kosong dengan prevalensi 4,0\% (\texttt{missing\_needs\_parking}) bernilai 4,88 simpangan baku saat bernilai 1, sehingga ketidakcocokannya menyumbang 23,8 pada jarak L2 kuadrat tetapi 4,88 pada jarak L1. Nilai ekstrem juga ada pada fitur numerik: nilai maksimum \texttt{prior\_completed\_bookings} berada 36 simpangan baku dari rata-rata dan \texttt{num\_children} 24,9 simpangan baku (data pengembangan). Pada L2 kuadrat sumbangan keduanya masing-masing sekitar 1.300 dan 620, pada L1 hanya 36 dan 25. Kedua, jarak L1 menjumlahkan sumbangan tiap fitur tanpa saling memperkuat, sehingga pembobotan per fitur (Subbab~\ref{subsubsec:svm-agop-teori}) bekerja langsung sebagai skala tiap sumbu. Kedua alasan ini adalah penjelasan yang masuk akal, bukan bukti. Buktinya adalah hasil empiris pada Subbab~\ref{subsec:svm-validasi} dan \ref{subsec:svm-hasil}. Karena bentuk eksponensial dari jarak L1 adalah kernel positif semidefinit, masalah dual \eqref{eq:svm-dual} tetap konkaf dan SMO konvergen. Implementasi \texttt{"laplacian"} tidak dapat memakai trik $\lVert a-b\rVert^2 = \lVert a\rVert^2 + \lVert b\rVert^2 - 2a\cdot b$ (perkalian matriks) seperti \textit{Gaussian}, sehingga jaraknya dihitung per kolom fitur (Kode~\ref{lst:svm-l1}) dan lebih lambat per baris kernel.

\subsubsection{Algoritma SMO}
Masalah dual \eqref{eq:svm-dual} adalah \textit{quadratic programming} dengan $n$ variabel, satu kendala persamaan, dan kendala kotak. \textit{Sequential Minimal Optimization} (SMO) \cite{platt1998smo} menyelesaikannya dengan mengubah dua variabel $\alpha_i, \alpha_j$ pada setiap langkah. Dua adalah jumlah paling sedikit yang memungkinkan kendala $\sum_k \alpha_k y_k = 0$ tetap terpenuhi, dan masalah dua variabel memiliki solusi tertutup, sehingga tidak diperlukan \textit{solver} QP. Implementasi mengikuti bentuk SMO pada LIBSVM \cite{chang2011libsvm} dengan pemilihan pasangan orde kedua \cite{fan2005wss}. Notasinya sebagai berikut.

\paragraph{Skor.} Untuk setiap titik $t$ didefinisikan
\begin{equation}
	s_t = y_t - \sum_{k}\alpha_k y_k K(x_k, x_t) = y_t - \big(f(x_t) - b\big),
	\label{eq:svm-score}
\end{equation}
yang pada kode adalah larik \texttt{score}. Pada awal pelatihan $\alpha = 0$ sehingga $s = y$. Gradien fungsi tujuan dual terhadap $\alpha_t$ sama dengan $y_t s_t$, sehingga $s_t$ cukup untuk memantau optimalitas. Selama pelatihan hanya vektor $s$, vektor $\alpha$, dan baris kernel yang sedang dipakai yang disimpan, bukan seluruh matriks kernel.

\paragraph{Himpunan indeks dan kondisi optimalitas.} Titik yang $\alpha$-nya masih boleh naik atau turun dikelompokkan menurut arah geraknya terhadap kendala $\sum\alpha_k y_k = 0$:
\begin{align}
	I_{\text{naik}} &= \{t : (y_t = +1,\ \alpha_t < C_t) \lor (y_t = -1,\ \alpha_t > 0)\}, \nonumber\\
	I_{\text{turun}} &= \{t : (y_t = +1,\ \alpha_t > 0) \lor (y_t = -1,\ \alpha_t < C_t)\},
	\label{eq:svm-indeks}
\end{align}
yang pada kode adalah \texttt{can\_rise} dan \texttt{can\_fall}. Menerapkan Persamaan~\eqref{eq:svm-kkt} pada $s_t = y_t - f(x_t) + b$ menunjukkan bahwa solusi optimum ada jika dan hanya jika terdapat $b$ dengan $s_t \le b$ untuk $t \in I_{\text{naik}}$ dan $s_t \ge b$ untuk $t \in I_{\text{turun}}$. Titik \textit{free support vector} memenuhi keduanya sehingga $s_t = b$. Dengan demikian optimalitas setara dengan
\begin{equation}
	m(\alpha) - M(\alpha) \le 0, \qquad m(\alpha) = \max_{t \in I_{\text{naik}}} s_t, \quad M(\alpha) = \min_{t \in I_{\text{turun}}} s_t .
\end{equation}
Pada praktiknya pelatihan berhenti jika $m(\alpha) - M(\alpha) < \tau$ dengan $\tau$ = \texttt{tol} (bawaan $10^{-3}$, sama dengan \textit{scikit-learn}), atau jika \texttt{n\_iter\_} mencapai \texttt{max\_iter}. Pada kasus kedua, \texttt{fit} mengeluarkan peringatan dan tetap mengembalikan model yang dapat dipakai.

\paragraph{Pemilihan pasangan orde kedua.} Titik pertama dipilih sebagai pelanggar terbesar, $i = \arg\max_{t \in I_{\text{naik}}} s_t$. Titik kedua dipilih dari $I_{\text{turun}}$ dengan menimbang kenaikan fungsi tujuan yang sebenarnya, bukan hanya selisih skor. Jika $\alpha_i y_i$ naik sebesar $\delta$ dan $\alpha_j y_j$ turun sebesar $\delta$ (kendala persamaan tetap terjaga), fungsi tujuan dual berubah menurut
\begin{equation}
	\Delta D(\delta) = \delta\,(s_i - s_j) - \tfrac{1}{2}\,a_{ij}\,\delta^2, \qquad a_{ij} = K_{ii} + K_{jj} - 2K_{ij},
\end{equation}
dengan $a_{ij} \ge 0$ kelengkungan di sepanjang arah tersebut. Parabola ini mencapai puncak di $\delta^{*} = (s_i - s_j)/a_{ij}$ dengan kenaikan $(s_i - s_j)^2 / (2a_{ij})$. Karena itu
\begin{equation}
	g_t = s_i - s_t, \qquad a_{it} = \max\big(K_{ii} + K_{tt} - 2K_{it},\ \tau_a\big), \qquad
	j = \arg\max_{t \in I_{\text{turun}},\ g_t > 0}\ \frac{g_t^{2}}{a_{it}},
	\label{eq:svm-pasangan}
\end{equation}
dengan $\tau_a = 10^{-12}$ (konstanta \texttt{TAU}, mengikuti LIBSVM) agar tidak terjadi pembagian dengan nol bila kelengkungan nol atau sedikit negatif akibat pembulatan. Bandingkan dengan pemilihan orde pertama yang hanya mengambil $t$ dengan $s_t$ terkecil: pasangan itu memiliki selisih skor terbesar, tetapi bisa memiliki kelengkungan besar sehingga langkahnya pendek dan kenaikan tujuannya kecil.

\paragraph{Langkah dan pemotongan ke kotak.} Langkah tanpa kendala adalah $\delta^{*}$. Langkah tersebut dipotong oleh sisa ruang di dalam kotak $[0, C_t]$ untuk kedua variabel:
\begin{equation}
	r_i = \begin{cases} C_i - \alpha_i & y_i = +1 \\ \alpha_i & y_i = -1 \end{cases}, \qquad
	r_j = \begin{cases} \alpha_j & y_j = +1 \\ C_j - \alpha_j & y_j = -1 \end{cases}, \qquad
	\delta = \min\{\delta^{*},\ r_i,\ r_j\},
\end{equation}
dan pembaruannya
\begin{equation}
	\alpha_i \leftarrow \alpha_i + y_i\,\delta, \qquad \alpha_j \leftarrow \alpha_j - y_j\,\delta, \qquad
	s_t \leftarrow s_t - \delta\,\big(K_{it} - K_{jt}\big)\quad \forall t .
	\label{eq:svm-update}
\end{equation}
Karena $y_i \cdot y_i\delta + y_j\cdot(-y_j\delta) = 0$, kendala $\sum\alpha_k y_k = 0$ tidak pernah dilanggar. Pembaruan skor muncul dari Persamaan~\eqref{eq:svm-score}: perubahan $\alpha_i y_i$ sebesar $+\delta$ dan $\alpha_j y_j$ sebesar $-\delta$ mengubah $f(x_t) - b$ sebesar $\delta(K_{it} - K_{jt})$. Jika langkah terpotong oleh batas (\texttt{step == room}), variabel yang bersangkutan langsung diset tepat ke $0$ atau $C_t$ oleh fungsi \texttt{\_move}, bukan hasil penjumlahan, agar galat pembulatan tidak membuat $\alpha$ berada sedikit di dalam atau di luar kotak. Setelah itu keanggotaan $i$ dan $j$ pada $I_{\text{naik}}$ dan $I_{\text{turun}}$ diperbarui.

\paragraph{Bias.} Titik \textit{free support vector} memenuhi $s_t = b$, sehingga
\begin{equation}
	b = \frac{1}{|F|}\sum_{t \in F} s_t, \qquad F = \{t : 0 < \alpha_t < C_t\},
\end{equation}
mengikuti LIBSVM \cite{chang2011libsvm}. Jika tidak ada \textit{free support vector}, bias diambil dari titik tengah $(m + M)/2$.

\subsubsection{Bobot Fitur dengan AGOP}
\label{subsubsec:svm-agop-teori}
Kernel Laplacian pada model terbaik tidak memakai fitur standar secara langsung, melainkan $K(a, b) = \exp\big(-\sum_k w_k|a_k - b_k| / L\big)$ dengan $L = 20$ (jadi $\gamma = 1/L = 0{,}05$) dan bobot $w_k \ge 0$ per fitur. Memakai bobot ini setara dengan mengalikan setiap kolom fitur dengan $w_k$ sebelum dimasukkan ke \texttt{SVM}, dan itulah yang dilakukan kode (\texttt{A * w}). Bobot dipelajari dari data latih dengan prosedur yang diadaptasi dari \textit{Average Gradient Outer Product} (AGOP) pada \textit{Recursive Feature Machine} \cite{radhakrishnan2024rfm}. Versi di sini disederhanakan: hanya diagonal AGOP yang dipakai, jaraknya L1, dan bobot adalah akar dari RMS gradien. Prosedur ini bukan replikasi penuh makalah tersebut. Satu putaran terdiri atas langkah berikut, dengan $X$ fitur terstandardisasi ($n \times d$) dan target $t_i = y_i - \tfrac12 \in \{\pm\tfrac12\}$. Pada putaran pertama $w = \mathbf{1}$, dan putaran dilakukan dua kali (\texttt{ROUNDS = 2}).

\begin{enumerate}
	\item \textbf{Model pengganti.} Dipilih $m = 1.500$ pusat acak $z_c$ dari baris latih (\texttt{CENTRES}), dan dipasang model $f(x) = \sum_{c=1}^{m} \beta_c\,K_w(x, z_c)$ dengan $K_w$ kernel Laplacian berbobot $w$ saat ini. Model ini adalah \textit{kernel ridge regression} \cite{saunders1998ridge} yang disederhanakan dengan pendekatan Nystr\"om \cite{williams2001nystrom}, yaitu pusat dibatasi pada himpunan bagian $m \ll n$ baris:
	\begin{equation}
		\min_{\beta}\ \sum_{i=1}^{n}\big(f(x_i) - t_i\big)^2 + \lambda\,\beta^{\top}K_{ZZ}\beta
		\ \Longrightarrow\
		\big(K_{XZ}^{\top}K_{XZ} + \lambda K_{ZZ} + \epsilon I\big)\beta = K_{XZ}^{\top}t,
		\label{eq:svm-nystrom}
	\end{equation}
	dengan $\lambda = 0{,}3$ (\texttt{RIDGE}) dan $\epsilon = 10^{-6}$ kali rata-rata diagonal matriks sistem (penstabil numerik). $K_{XZ}^{\top}K_{XZ}$ dijumlahkan per blok 2.000 baris agar memori tidak bergantung pada $n$.
	\item \textbf{Gradien.} Turunan model terhadap fitur ke-$k$ pada titik $x$ (turunan $|u|$ adalah $\operatorname{sign}(u)$, terdefinisi hampir di mana-mana) adalah
	\begin{equation}
		\frac{\partial f}{\partial x_k}(x) = -\frac{w_k}{L}\sum_{c=1}^{m}\beta_c\,K_w(x, z_c)\,\operatorname{sign}\big(x_k - z_{ck}\big).
		\label{eq:svm-gradien}
	\end{equation}
	Gradien dihitung pada $q = 4.000$ baris acak (\texttt{GRADIENT\_ROWS}). Pada kode, jumlahnya dihitung pada koordinat yang sudah berbobot, lalu dikalikan $w_k$ sesuai aturan rantai.
	\item \textbf{Diagonal AGOP dan bobot baru.} Rata-rata kuadrat gradien per fitur adalah diagonal AGOP,
	\begin{equation}
		G_k = \frac{1}{q}\sum_{r=1}^{q}\Big(\frac{\partial f}{\partial x_k}(x_r)\Big)^{2}, \qquad
		w_k^{\text{baru}} \propto \sqrt{\sqrt{G_k}} = G_k^{1/4}, \qquad \sum_{k=1}^{d} \big(w_k^{\text{baru}}\big)^2 = d .
		\label{eq:svm-bobot}
	\end{equation}
	Fitur yang membuat prediksi model pengganti berubah banyak ketika digeser mendapat bobot besar. Pangkat $1/4$ pada $G_k$ (akar dari RMS gradien) dipilih lewat validasi silang (pangkat 0,35 dan 0,7 pada RMS lebih rendah, Tabel~\ref{tab:svm-ablasi}). Normalisasi $\sum_k w_k^2 = d$ menjaga rata-rata kuadrat bobot tetap 1, sehingga skala jarak keseluruhan tidak berubah dan $L = 20$ tetap bermakna. Bobot akhir dibulatkan ke tiga desimal.
\end{enumerate}

Bobot ini bekerja sebagai pemilihan fitur yang halus: fitur tidak relevan tidak dibuang tetapi diperkecil pengaruhnya pada jarak (Tabel~\ref{tab:svm-bobot}). Prosedur ini jauh lebih mahal daripada standardisasi. Pada $n = 27.309$, $m = 1.500$, dan $d = 46$, pembentukan $K_{XZ}^{\top}K_{XZ}$ membutuhkan sekitar $n m^2 \approx 6{,}1 \times 10^{10}$ operasi kali-tambah (BLAS), jarak L1 antara seluruh baris dan pusat membutuhkan $n m d \approx 1{,}9 \times 10^{9}$ operasi, dan penyelesaian sistem $m \times m$ sekitar $m^3 \approx 3{,}4 \times 10^{9}$ operasi, semuanya per putaran. Bobot dihitung ulang dari awal untuk setiap pembagian data, sehingga tidak bocor dari data validasi (Subbab~\ref{subsec:svm-preprocessing}).

\subsubsection{Struktur Kelas dan Fungsi}
Kelas \texttt{SVM} berada pada \nolinkurl{src/algoritma/svm/svm.py} dan hanya memakai NumPy. Fitur dibangun oleh fungsi pada \nolinkurl{src/kaggle/booking_features.py} (pandas dan NumPy). Tabel~\ref{tab:svm-fungsi} merangkum fungsi dan metode pada \texttt{svm.py}, dan Tabel~\ref{tab:svm-fitur-fungsi} merangkum fungsi pada \texttt{booking\_features.py}.

\begin{table}[H]
	\centering
	\caption{Deskripsi kelas dan fungsi SVM (\texttt{src/algoritma/svm/svm.py})}
	\label{tab:svm-fungsi}
	\small
	\begin{tabularx}{\textwidth}{l Y}
		\toprule
		\textbf{Fungsi/Metode} & \textbf{Deskripsi} \\
		\midrule
		\texttt{KERNELS}, \texttt{TAU} & Tuple nama kernel yang sah (\texttt{linear}, \texttt{poly}, \texttt{rbf}, \texttt{laplacian}) dan batas bawah kelengkungan $\tau_a = 10^{-12}$. \\
		\texttt{PREDICT\_BATCH\_ROWS}, \texttt{LAPLACIAN\_BATCH\_ROWS} & Jumlah baris data uji yang diproses sekaligus oleh \texttt{decision\_function} (4.096; 256 untuk Laplacian) agar memori matriks kernel terbatas. \\
		\texttt{\_l1\_distances(A, B)} & Matriks jarak L1 antara baris \texttt{A} dan \texttt{B}, dijumlahkan per kolom fitur agar tidak membentuk larik tiga dimensi. Kasus satu baris dihitung langsung. \\
		\texttt{\_KernelRowCache.\_\_init\_\_} & Menyimpan fungsi penghitung baris kernel dan kapasitas \textit{cache} $\max\big(2, \lfloor \texttt{cache\_mb}\cdot 10^6 / (8n)\rfloor\big)$ baris. \\
		\texttt{\_KernelRowCache.\_\_getitem\_\_} & Mengembalikan baris kernel ke-$i$ dari \textit{cache} (kebijakan \textit{least recently used} dengan \texttt{OrderedDict}), atau menghitungnya dan membuang baris terlama jika \textit{cache} penuh. \\
		\texttt{SVM.\_\_init\_\_(...)} & Menyimpan hiperparameter: \texttt{C}, \texttt{kernel}, \texttt{degree}, \texttt{gamma}, \texttt{coef0}, \texttt{tol}, \texttt{max\_iter}, \texttt{class\_weight}, dan \texttt{cache\_mb}. Kernel yang tidak dikenal ditolak dengan \texttt{ValueError}. \\
		\texttt{fit(X, y)} & Memvalidasi masukan (dua kelas, tanpa \texttt{NaN}), memetakan label ke $\pm 1$, menghitung batas kotak $C_i$ dan $\gamma$, memanggil \texttt{\_solve}, lalu menyimpan \texttt{support\_}, \texttt{support\_vectors\_}, \texttt{dual\_coef\_} ($\alpha_i y_i$), \texttt{intercept\_}, dan \texttt{coef\_} (hanya kernel \textit{linear}). \\
		\texttt{decision\_function(X)} & Skor $f(x)$ menurut Persamaan~\eqref{eq:svm-decision}, dihitung per kelompok baris. \\
		\texttt{predict(X, threshold)} & Mengembalikan label kelas positif jika skor lebih besar dari \texttt{threshold} (bawaan 0). Label asli (bukan hanya $\pm 1$) dikembalikan. \\
		\texttt{save(path)} & Menyimpan seluruh objek model dengan \texttt{pickle}. \\
		\texttt{load(path)} & \textit{Class method} yang memuat model dari berkas \textit{pickle}. \\
		\texttt{\_solve(X, sign, upper)} & Inti SMO: memilih pasangan, memperbarui $\alpha$ dan skor, memeriksa kriteria berhenti, dan menghitung bias. Mengembalikan $(\alpha, b)$ dan mengisi \texttt{n\_iter\_}. \\
		\texttt{\_move(value, delta, hits\_bound, upper)} & Memperbarui satu $\alpha$; jika langkah terpotong batas, nilainya diset persis ke $0$ atau \texttt{upper}. \\
		\texttt{\_gram(A, B, B\_squared\_norms)} & Matriks kernel antara baris \texttt{A} dan \texttt{B} untuk kernel yang dipilih. \\
		\texttt{\_gram\_diagonal(squared\_norms)} & Diagonal matriks kernel, $K_{ii}$, tanpa membentuk seluruh matriks. \\
		\texttt{\_resolve\_gamma(X)} & Mengubah \texttt{"scale"} dan \texttt{"auto"} menjadi nilai $\gamma$ numerik. \\
		\texttt{\_box\_limits(y)} & Batas atas $C_i$ tiap titik dari \texttt{class\_weight} (\texttt{None}, \texttt{"balanced"}, atau kamus). \\
		\texttt{\_check\_fitted()} & Melempar \texttt{RuntimeError} jika \texttt{predict} dipanggil sebelum \texttt{fit}. \\
		\texttt{\_as\_matrix(X)} & Mengubah masukan menjadi larik \texttt{float64} \textit{contiguous} dan menolak \texttt{NaN} atau tak hingga. \\
		\bottomrule
	\end{tabularx}
\end{table}

\begin{table}[H]
	\centering
	\caption{Fungsi pembangun fitur dan bobot (\texttt{src/kaggle/booking\_features.py})}
	\label{tab:svm-fitur-fungsi}
	\small
	\begin{tabularx}{\textwidth}{l Y}
		\toprule
		\textbf{Fungsi} & \textbf{Deskripsi} \\
		\midrule
		\texttt{load\_data(data\_dir)} & Membaca \texttt{train.csv} dan \texttt{test.csv} dengan \texttt{id} sebagai teks. \\
		\texttt{clean(train)} & Membuang kolom \texttt{id} lalu baris duplikat (data latih saja). \\
		\texttt{add\_booking\_dates(df)} & Menambahkan \texttt{arr\_ord} (hari sejak 2023-01-01), \texttt{book\_ord} (\texttt{arr\_ord} dikurangi \textit{lead time}), dan \texttt{arr\_dow} (hari dalam minggu). \\
		\texttt{impute\_and\_flag(train, other)} & Imputasi median atau modus yang dipelajari dari data latih, dan penanda \texttt{missing\_*} untuk setiap kolom yang punya nilai kosong di data latih. \\
		\texttt{sibling\_pairs(reference, query, same)} & Mencari saudara tiap baris \texttt{query} pada \texttt{reference} (kunci keras, toleransi, kolom lunak) beserta bobot kedekatan \textit{lead time}. \\
		\texttt{shared\_neighbour\_weights(...)} & Mengalikan bobot pasangan dengan akar indeks Jaccard antar-himpunan saudara. \\
		\texttt{build\_features(train, labels, other, shared\_neighbours)} & Membangun seluruh fitur: tanggal, imputasi, fitur saudara (\texttt{sib\_te}, \texttt{sib\_n}) secara \textit{out-of-fold}, lama menginap, dan \textit{one-hot}. Kolom \texttt{other} disejajarkan dengan kolom data latih. \\
		\texttt{standardize(train, other)} & Transformasi akar, logit, dan $\log(1+x)$ pada tiga kolom, lalu \textit{z-score}; kolom dengan simpangan baku nol dibuang. Mengembalikan nama kolom. \\
		\texttt{laplacian\_gram(A, B, scale)} & Matriks kernel Laplacian $\exp(-\lVert a-b\rVert_1/L)$ untuk penghitungan bobot fitur. \\
		\texttt{feature\_weights(X, y, scale)} & Bobot fitur AGOP diagonal (Persamaan~\eqref{eq:svm-nystrom}--\eqref{eq:svm-bobot}). \\
		\bottomrule
	\end{tabularx}
\end{table}

Konstanta yang dipakai pembangun fitur dan bobot ditunjukkan pada Tabel~\ref{tab:svm-konstanta}.

\begin{table}[H]
	\centering
	\caption{Konstanta pada \texttt{booking\_features.py}}
	\label{tab:svm-konstanta}
	\small
	\begin{tabularx}{\textwidth}{l c Y}
		\toprule
		\textbf{Konstanta} & \textbf{Nilai} & \textbf{Arti} \\
		\midrule
		\texttt{LEAD\_TOLERANCE} & 10 & Selisih \textit{lead time} paling besar antara dua saudara (hari). \\
		\texttt{RATE\_TOLERANCE} & 0,05 & Selisih tarif paling besar, sebagai pecahan dari tarif yang lebih besar. \\
		\texttt{MISSING\_LEAD\_WEIGHT} & 0,1 & Bobot pasangan jika salah satu \textit{lead time} kosong. \\
		\texttt{INNER\_FOLDS} & 5 & Jumlah kelompok acak untuk fitur saudara \textit{out-of-fold}. \\
		\texttt{CENTRES} & 1.500 & Jumlah pusat Nystr\"om untuk bobot fitur. \\
		\texttt{GRADIENT\_ROWS} & 4.000 & Jumlah baris untuk menghitung gradien. \\
		\texttt{RIDGE} & 0,3 & Kekuatan regularisasi $\lambda$ pada Persamaan~\eqref{eq:svm-nystrom}. \\
		\texttt{ROUNDS} & 2 & Jumlah putaran pembaruan bobot. \\
		\bottomrule
	\end{tabularx}
\end{table}

\subsubsection{\textit{Source Code}}
Nomor baris yang disebut pada teks adalah nomor baris di dalam masing-masing potongan kode (dimulai dari 1), bukan nomor baris pada berkas aslinya.

Kode~\ref{lst:svm-fit} adalah \texttt{fit}. Label dipetakan ke $\pm 1$ dengan kelas berlabel lebih besar (\texttt{classes\_[1]}) sebagai kelas positif, sehingga skor positif berarti kelas batal. Hanya titik dengan $\alpha > 0$ yang disimpan sebagai \textit{support vector}, dan \texttt{dual\_coef\_} menyimpan $\alpha_i y_i$ yang langsung dipakai pada Persamaan~\eqref{eq:svm-decision}.

\kode{../../src/algoritma/svm/svm.py}{59}{78}{Metode \texttt{fit}: validasi, pemetaan label, SMO, dan penyimpanan \textit{support vector}}{lst:svm-fit}

Kode~\ref{lst:svm-l1} dan~\ref{lst:svm-gram} menghitung matriks kernel. Jarak L1 dihitung per kolom fitur: untuk $d = 46$ kolom, satu matriks berukuran $|A| \times |B|$ ditambahkan per kolom, sehingga memori hanya $O(|A||B|)$ dan bukan $O(|A||B|d)$. Kernel \textit{linear}, \textit{polynomial}, dan \textit{Gaussian} memakai perkalian matriks \texttt{A @ B.T} yang dioptimalkan BLAS; pada \textit{Gaussian} jarak kuadrat dihitung dari $\lVert a\rVert^2 + \lVert b\rVert^2 - 2a\cdot b$ lalu dipotong ke nol untuk menghindari nilai negatif kecil akibat pembulatan.

\kode{../../src/algoritma/svm/svm.py}{16}{22}{Jarak L1 per kolom fitur}{lst:svm-l1}

\kode{../../src/algoritma/svm/svm.py}{177}{195}{Matriks kernel dan diagonalnya untuk keempat kernel}{lst:svm-gram}

Kode~\ref{lst:svm-smo-a} adalah awal SMO. Larik \texttt{score}, \texttt{can\_rise}, dan \texttt{can\_fall} berturut-turut adalah $s$, $I_{\text{naik}}$, dan $I_{\text{turun}}$ pada Persamaan~\eqref{eq:svm-score} dan \eqref{eq:svm-indeks}. Pada $\alpha = 0$, \texttt{can\_rise} adalah titik berlabel $+1$ dan \texttt{can\_fall} berlabel $-1$. Setiap iterasi memilih $i$ (baris~10), menghentikan pelatihan jika $m - M < \tau$ (baris~13), lalu memilih $j$ dengan kenaikan $g_t^2/a_{it}$ terbesar (baris~19--23) sesuai Persamaan~\eqref{eq:svm-pasangan}. Baris kernel $i$ diambil dari \textit{cache}.

\kode{../../src/algoritma/svm/svm.py}{117}{139}{SMO: inisialisasi, kriteria berhenti, dan pemilihan pasangan orde kedua}{lst:svm-smo-a}

Kode~\ref{lst:svm-smo-b} menerapkan Persamaan~\eqref{eq:svm-update}: langkah dipotong ke ruang yang tersisa dalam kotak (baris~3--6), kedua $\alpha$ digerakkan ke arah yang menjaga $\sum\alpha_k y_k = 0$, skor seluruh titik diperbarui dengan selisih dua baris kernel (baris~10), dan keanggotaan himpunan indeks kedua titik dihitung ulang. Per iterasi hanya dua baris kernel yang diperlukan, dan seluruh operasi lain berupa operasi vektor sepanjang $n$.

\kode{../../src/algoritma/svm/svm.py}{141}{159}{SMO: langkah, pemotongan ke kotak, dan pembaruan skor}{lst:svm-smo-b}

Kode~\ref{lst:svm-smo-c} menangani kasus \texttt{max\_iter} tercapai sebelum konvergen (peringatan, model tetap dipakai) lalu menghitung bias dari rata-rata \textit{free support vector}, dan memuat fungsi \texttt{\_move}.

\kode{../../src/algoritma/svm/svm.py}{161}{175}{SMO: peringatan, bias, dan \texttt{\_move}}{lst:svm-smo-c}

Kode~\ref{lst:svm-agop} adalah penghitungan bobot fitur. Setiap putaran membentuk $K_{XZ}^{\top}K_{XZ}$ dan $K_{XZ}^{\top}t$ per blok (baris~10--13), menyelesaikan sistem Nystr\"om (baris~14--16), menghitung gradien Persamaan~\eqref{eq:svm-gradien} per fitur (baris~19--23), lalu memperbarui bobot menurut Persamaan~\eqref{eq:svm-bobot} (baris~24--25). Dua generator acak ber-\textit{seed} tetap (0 untuk pusat dan 1 untuk baris gradien) membuat hasilnya \textit{reproducible}.

\kode{../../src/kaggle/booking_features.py}{174}{199}{Bobot fitur AGOP diagonal dengan model pengganti Nystr\"om}{lst:svm-agop}

\subsubsection{Kompleksitas dan \textit{Cache}}
\paragraph{Waktu.} Satu iterasi SMO mengerjakan beberapa operasi vektor sepanjang $n$ (pencarian $i$, $j$, dan pembaruan skor), sehingga $O(n)$, ditambah penghitungan baris kernel jika tidak ada di \textit{cache}, yaitu $O(nd)$ per baris. Total waktu adalah jumlah iterasi dikali biaya itu. Jumlah iterasi tidak memiliki batas yang rapat: pada model akhir ($n = 27.309$) SMO membutuhkan 22.302 iterasi (0,82$n$), sedangkan kernel \textit{linear} pada data yang sama dengan $C = 8$ membutuhkan 936.187 iterasi untuk hanya 6.000 baris latih (Subbab~\ref{subsec:svm-hasil}). Implementasi tidak memakai \textit{shrinking}, yaitu membuang sementara titik yang sudah pasti berada di batas, sehingga setiap iterasi tetap memeriksa seluruh $n$ titik.

\paragraph{Memori dan \textit{cache}.} Menyimpan seluruh matriks kernel membutuhkan $8n^2$ byte: 5,97 GB untuk $n = 27.309$, 3,82 GB untuk data pengembangan ($n = 21.847$), dan 288 MB untuk 6.000 baris. Karena itu hanya baris yang dipakai yang disimpan dalam \textit{cache} \textit{least recently used} dengan kapasitas $\lfloor \texttt{cache\_mb}\cdot 10^6 / (8n)\rfloor$ baris. Pada model akhir (\texttt{cache\_mb=400}, $n = 27.309$), kapasitasnya 1.830 baris, atau 6,7\% dari seluruh baris. \textit{Cache} hanya memengaruhi kecepatan, bukan hasil: baris yang dihitung ulang identik dengan baris yang disimpan, dan pengujian memastikan \textit{cache} sekecil dua baris menghasilkan \texttt{n\_iter\_} dan skor yang sama persis dengan \textit{cache} besar (Subbab~\ref{subsubsec:svm-optimasi} memuat pengukuran waktunya).

\paragraph{Prediksi.} \texttt{decision\_function} memproses data uji per kelompok (4.096 baris, atau 256 untuk Laplacian), sehingga matriks kernel sementara paling besar $256 \times 8.662$ elemen (sekitar 17,7 MB) pada model akhir. Biaya prediksi $O(m\,n_{sv}\,d)$ untuk $m$ baris uji dan $n_{sv}$ \textit{support vector}.

\subsubsection{\textit{Save} dan \textit{Load} Model}
Model disimpan dengan \texttt{pickle}: \texttt{save} menulis seluruh objek \texttt{SVM}, dan \texttt{load} membacanya kembali. Berkas berisi \textit{support vector} (\texttt{support\_vectors\_}), koefisien dual (\texttt{dual\_coef\_}), bias (\texttt{intercept\_}), $\gamma$ yang sudah dihitung (\texttt{gamma\_}), label kelas (\texttt{classes\_}), indeks \textit{support vector} (\texttt{support\_}), jumlah iterasi (\texttt{n\_iter\_}), dan hiperparameter. Model akhir (8.662 \textit{support vector}, 46 fitur) berukuran 3,3 MB di \nolinkurl{output/svm/svm_model.pkl}. \textit{Pickle} dipilih karena objek memuat larik dan skalar sehingga tidak perlu format khusus. Kelemahannya, \texttt{pickle.load} dapat menjalankan kode sembarang dari berkas yang dimuat, sehingga hanya berkas dari sumber tepercaya yang boleh dimuat (pada LR dipakai \texttt{.npz} yang tidak memiliki risiko ini, Subbab~\ref{subsec:lr-implementasi}).

\begin{lstlisting}[caption={Contoh penyimpanan dan pemuatan model},label={lst:svm-saveload}]
model = SVM(kernel="laplacian", gamma=1 / 20, C=2.0, tol=1e-3).fit(A * w, y)
model.save("output/svm/svm_model.pkl")
loaded = SVM.load("output/svm/svm_model.pkl")
assert np.array_equal(loaded.decision_function(B * w), model.decision_function(B * w))
\end{lstlisting}

Notebook memeriksa bahwa nilai keputusan model hasil \texttt{load} identik secara bit dengan model aslinya pada model \textit{hold-out}, dan model \textit{submission} dimuat kembali dari berkas sebelum dipakai memprediksi. Ada keterbatasan yang sama seperti LR: berkas hanya menyimpan model, bukan pembangun fitur. \textit{Support vector} tersimpan dalam ruang fitur yang sudah dikalikan bobot $w$, tetapi $w$, statistik standardisasi, dan tabel saudara tidak ikut tersimpan. Untuk memprediksi data baru, \texttt{build\_features}, \texttt{standardize}, dan \texttt{feature\_weights} harus dijalankan ulang pada data latih yang sama (deterministik, sehingga hasilnya identik) lalu fitur data baru dikalikan $w$. Menyimpan $w$, rata-rata, dan simpangan baku ke berkas yang sama adalah perbaikan yang wajar, tetapi tabel saudara tetap memerlukan data latih karena fitur \texttt{sib\_te} mencari pesanan serupa di dalamnya.

\subsubsection{Pengujian}
\label{subsubsec:svm-uji}
Pengujian ada pada \nolinkurl{test/test_svm.py}, berisi 22 metode uji (dijalankan dengan \texttt{python}\ \nolinkurl{test/test_svm.py}), semuanya lolos dalam sekitar 0,7 detik. Pengujian dibagi menjadi tujuh kelompok.
\begin{itemize}
	\item \textbf{Data terpisah (2).} Kernel \textit{linear} mengklasifikasikan dua gugus terpisah dengan akurasi 100\%, dan kernel \textit{Gaussian} menyelesaikan data \texttt{make\_moons} (akurasi $> 0{,}99$) yang tidak dapat diselesaikan kernel \textit{linear} ($< 0{,}95$).
	\item \textbf{Kesesuaian dengan LIBSVM (5).} Pada 8 konfigurasi (\textit{linear} dengan dua nilai $C$, dua konfigurasi \textit{polynomial}, dua konfigurasi \textit{Gaussian}, dan dua konfigurasi \textit{Gaussian} dengan \texttt{class\_weight}) dengan \texttt{tol} $10^{-5}$, nilai keputusan pada data baru berbeda paling besar $5 \times 10^{-3}$ dari \texttt{SVC}, dan jumlah \textit{support vector} berbeda paling banyak 2. Kernel Laplacian dibandingkan dengan \texttt{SVC} yang memakai fungsi kernel dari \texttt{sklearn.metrics.pairwise.laplacian\_kernel} (batas yang sama). Dua uji lain memeriksa \texttt{gamma="scale"} dan \texttt{gamma="auto"} terhadap rumusnya.
	\item \textbf{Kendala dual (3).} $\sum_i \alpha_i y_i = 0$ (galat paling besar $10^{-8}$), kotak batas dengan bobot kelas (positif $\le 6$, negatif $\ge -2$ untuk $C = 2$, bobot 3:1), dan \texttt{coef\_} \textit{linear} yang mereproduksi \texttt{decision\_function}.
	\item \textbf{\textit{Cache} (1).} \textit{Cache} sekecil 0,01 MB memberi \texttt{n\_iter\_} dan skor identik dengan \textit{cache} 100 MB.
	\item \textbf{Label (3).} Label asli (mis. teks) dikembalikan, kelas dengan urutan lebih besar adalah kelas positif, dan \texttt{threshold} yang lebih rendah menambah prediksi positif.
	\item \textbf{Persistensi (1).} Model hasil \texttt{load} memprediksi identik.
	\item \textbf{Validasi masukan (7).} Kernel tidak dikenal, satu kelas, lebih dari dua kelas, \texttt{NaN}, \texttt{predict} sebelum \texttt{fit}, \texttt{max\_iter} tercapai (ada peringatan, model tetap dapat dipakai), dan pelatihan yang konvergen tidak mengeluarkan peringatan.
\end{itemize}
Pengujian ini memeriksa kebenaran pada data buatan berukuran 400 baris. Kesetaraan pada data sebenarnya dan pada kernel Laplacian dengan bobot fitur diperiksa pada Subbab~\ref{subsec:svm-hasil}.

\subsection{\textit{Data Cleaning} dan \textit{Preprocessing}}
\label{subsec:svm-preprocessing}

\todo{Jelaskan \textit{pipeline} khusus SVM beserta alasan tiap langkah, dikaitkan dengan temuan EDA (Subbab~\ref{sec:eda}) dan karakteristik algoritma. Semua transformasi di-\textit{fit} pada data latih saja lalu diterapkan ke data uji \cite{slide-dataprep}.}

\subsubsection{Penanganan \textit{Missing Values}}
\todo{SVM tidak menerima \textit{null}, jadi imputasi wajib. Jelaskan strategi per kolom dan alasannya.}

\subsubsection{Data Duplikat dan Data Tidak Valid}
\todo{Jumlah duplikat, nilai tidak wajar, atau \textit{outlier} yang ditangani untuk SVM dan caranya. Baris data uji tidak boleh dibuang.}

\subsubsection{Konstruksi Fitur dan \textit{Encoding}}
\todo{SVM butuh fitur numerik. Jelaskan \textit{encoding} kategorikal dan pengaruhnya ke kernel.}

\subsubsection{Normalisasi dan Standarisasi}
\todo{Wajib, terutama untuk kernel RBF dan \textit{polynomial} yang sensitif terhadap skala.}

\subsubsection{Penanganan \textit{Imbalanced Data}}
\todo{Misalnya \textit{class weight} (nilai $C$ berbeda per kelas) atau penyesuaian \textit{threshold} pada \texttt{decision\_function}.}

\subsubsection{\textit{Feature Selection}}
\todo{Jelaskan seleksi fitur yang dipakai dan dampaknya ke waktu pelatihan. Sebutkan fitur akhir yang dipakai.}

\subsubsection{\textit{Dimensionality Reduction}}
\todo{Jelaskan apakah PCA dipakai (juga untuk visualisasi 2D bonus) dan dampaknya.}

\subsection{\textit{Modeling} dan Validasi}
\label{subsec:svm-validasi}

\subsubsection{Skema Validasi}
\todo{Skema validasi untuk SVM (\textit{stratified k-fold} atau \textit{train-test split}), nilai $k$/rasio, dan \textit{seed}. Jika sama dengan kerangka umum (Subbab~\ref{sec:kerangka-umum}), cukup rujuk dan jelaskan perbedaannya saja, misalnya \textit{subsampling} karena waktu komputasi.}

\subsubsection{\textit{Hyperparameter Tuning}}
\todo{Ruang pencarian (kernel, $C$, $\gamma$, \textit{degree}, \texttt{coef0}, ukuran sampel latih (karena kompleksitas $O(n^2)$)), metode pencarian, dan hasilnya. Sertakan tabel atau grafik skor CV terhadap hiperparameter.}

\begin{table}[H]
	\centering
	\caption{Hasil \textit{tuning} hiperparameter SVM}
	\label{tab:svm-tuning}
	\begin{tabular}{l l c}
		\toprule
		\textbf{Hiperparameter} & \textbf{Ruang pencarian} & \textbf{Terbaik} \\
		\midrule
		\todo{} & & \\
		\bottomrule
	\end{tabular}
\end{table}

\subsubsection{Model Final}
\todo{Konfigurasi final SVM, skor CV (rata-rata $\pm$ simpangan baku), dan lokasi model tersimpan di \texttt{output/}.}

\subsection{Hasil dan Perbandingan dengan \textit{Scikit-learn}}
\label{subsec:svm-hasil}

\todo{Bandingkan dengan \texttt{sklearn.svm.SVC} untuk tiap kernel dengan $C$, $\gamma$, dan \textit{degree} yang sama. Pakai \textit{preprocessing} dan \textit{fold} yang sama agar perbedaan hanya berasal dari implementasi \cite{pedregosa2011sklearn}.}

\begin{table}[H]
	\centering
	\caption{Perbandingan SVM \textit{from scratch} dan \textit{scikit-learn}}
	\label{tab:svm-hasil}
	\begin{tabular}{l c c c c c}
		\toprule
		\textbf{Implementasi} & \textbf{Macro F1} & \textbf{Accuracy} & \textbf{Precision} & \textbf{Recall} & \textbf{Waktu (s)} \\
		\midrule
		\textit{Scratch} -- \textit{Linear} & & & & & \\
		\textit{Sklearn} -- \textit{Linear} & & & & & \\
		\textit{Scratch} -- \textit{Polynomial} & & & & & \\
		\textit{Sklearn} -- \textit{Polynomial} & & & & & \\
		\textit{Scratch} -- \textit{Gaussian} & & & & & \\
		\textit{Sklearn} -- \textit{Gaussian} & & & & & \\
		\bottomrule
	\end{tabular}
\end{table}

\begin{figure}[H]
	\centering
	\placeholdergambar{\textit{Confusion matrix} \textit{from scratch} vs \textit{scikit-learn}}
	\caption{\textit{Confusion matrix} SVM}
	\label{fig:svm-cm}
\end{figure}


\subsubsection{Analisis}
\todo{Bandingkan juga jumlah \textit{support vector} dan nilai \texttt{decision\_function}. Jelaskan selisih akibat toleransi, pemilihan \textit{working set}, dan \textit{shrinking} pada LIBSVM.}

\subsection{Bonus SVM}
\label{subsec:svm-bonus}

\todo{Hapus subbab bonus yang tidak dikerjakan.}

\subsubsection{Visualisasi Bidang Pemisah}
\todo{Gambar bidang pemisah, \textit{margin}, dan \textit{support vector} untuk tiap kernel pada dua fitur atau hasil reduksi ke 2D. Jelaskan fitur yang dipilih dan tampilkan tiga gambarnya.}

\begin{figure}[H]
	\centering
	\placeholdergambar{Visualisasi Bidang Pemisah}
	\caption{Visualisasi Bidang Pemisah SVM}
	\label{fig:svm-bonus-0}
\end{figure}

\subsubsection{Algoritma Optimasi di Luar Materi Kelas}
\todo{Jika ada: jelaskan algoritmanya, mengapa dapat mengoptimasi model, bukti perbaikannya (sebelum vs sesudah), dan \textbf{referensinya (tanpa referensi bonus dianulir)}. Misalnya SMO dengan pemilihan \textit{working set} orde kedua dan \textit{kernel cache} \cite{platt1998smo,chang2011libsvm}.}

